\section{13. Conclusión}
Enlace Broadcast Platform ha evolucionado desde una necesidad específica de distribución de EJTV hacia una plataforma con capacidades reales de recepción, administración, distribución y supervisión de contenido audiovisual.

La existencia de señales principales como ENLACE y EJTV, la distribución hacia Cable Centro como primer Reader real, las pruebas técnicas realizadas con IMPACT, la identificación de las conexiones al servidor y la implementación del NOC demuestran que el proyecto ya dispone de componentes funcionales y evidencia operacional.

La evolución prevista ampliará esta supervisión hacia el análisis directo del contenido audiovisual, con detección de condiciones como ausencia de audio o video, congelamiento de imagen, pantalla negra y problemas asociados a la red, junto con mecanismos de correlación que ayuden a determinar el punto probable de origen de una incidencia. A esto se sumará, en una etapa futura, un multiviewer para la observación simultánea y en tiempo real de las señales entrantes.

El trabajo pendiente se concentra en consolidar las capacidades existentes, ampliar las pruebas físicas, incorporar nuevos receptores y fuentes de forma controlada, continuar la evolución del acceso web al NOC, desarrollar progresivamente el análisis audiovisual y preparar las siguientes etapas de multiviewer, escalamiento y transcodificación.

\begin{callout}Estado general: EBP cuenta con una base operacional funcional y una arquitectura de supervisión en evolución, orientada a integrar conectividad, estado de infraestructura, Health de protocolos, análisis audiovisual, trazabilidad de Readers y monitoreo visual en una misma plataforma.\end{callout}
